\documentclass{report}
\usepackage{amsmath,amssymb}
\setlength{\parindent}{0mm}
\setlength{\parskip}{1em}
\begin{document}
\begin{center}
\rule{6in}{1pt} \
{\large Kuo Liu \\
{\bf Hybrid First-order System Least Squares Finite Element Methods With Application to Stokes Equations}}

1320 Grandview Ave \\ Boulder \\ CO 80302
\\
{\tt kuol@colorado.edu}\\
Thomas Manteuffel\\
Stephen McCormick\\
John Ruge\\
Lei Tang\end{center}

\newcommand{\BB}{{\bf B}}
\newcommand{\BF}{{\bf F}}
\newcommand{\BL}{{\bf L}}
\newcommand{\BU}{{\bf U}}
\newcommand{\BV}{{\bf V}}
\newcommand{\BW}{{\bf W}}
\newcommand{\BZ}{{\bf Z}}

\newcommand{\bb}{{\bf b}}
\newcommand{\bc}{{\bf c}}
\newcommand{\bfe}{{\bf e}}
\newcommand{\bff}{{\bf f}}
\newcommand{\bg}{{\bf g}}
\newcommand{\bn}{{\bf n}}
\newcommand{\bp}{{\bf p}}
\newcommand{\br}{{\bf r}}
\newcommand{\bu}{{\bf u}}
\newcommand{\bv}{{\bf v}}
\newcommand{\bw}{{\bf w}}
\newcommand{\bx}{{\bf x}}
\newcommand{\by}{{\bf y}}
\newcommand{\bz}{{\bf z}}


\newcommand{\btau}{\mathbf{\tau}}
\newcommand{\bzero}{{\bf 0}}
\newcommand{\Bzero}{{\bf O}}
\newcommand{\bomega}{\mathbf{\omega}}

\newcommand{\CB}{{\cal B}}
\newcommand{\CC}{{\cal C}}
\newcommand{\CD}{{\cal D}}
\newcommand{\CF}{{\cal F}}
\newcommand{\CG}{{\cal G}}
\newcommand{\CH}{{\cal H}}
\newcommand{\CI}{{\cal I}}
\newcommand{\CL}{{\cal L}}
\newcommand{\CN}{{\cal N}}
\newcommand{\CO}{{\cal O}}
\newcommand{\CR}{{\cal R}}
\newcommand{\CT}{{\cal T}}
\newcommand{\CU}{{\cal U}}
\newcommand{\CV}{{\cal V}}
\newcommand{\CW}{{\cal W}}




In this talk, we combine the FOSLS method with the FOSLL$^*$ method to
create a Hybrid method.
The FOSLS approach minimizes the error, $\bfe^h = \bu^h - \bu$, over a finite element
subspace, $\CV^h$,
in the operator norm, $\min_{\bu^h\in\CV^h}\| L (\bu^h-\bu) \|$. The FOSLL$^*$ method
looks for an approximation
in the range of $L^*$, setting $\bu^h = L^*\bw^h$ and choosing $\bw^h \in \CW^h$, a
standard finite element space.
FOSLL$^*$ minimizes the $\BL^2$ norm of the error over $L^*(\CW^h)$, that is,
$\min_{\bw^h\in\CW^h} \| L^*\bw^h - \bu\|$.
FOSLS enjoys a locally sharp, globally reliable, and easily computable a posterior error
estimate, while FOSLL$^*$
does not.

The hybrid method attempts to retain the best properties of both FOSLS and FOSLL$^*$.
This is accomplished by combining the FOSLS functional, the FOSLL$^*$ functional, and an
intermediate term that draws them together. The Hybrid method produces an approximation,
$\bu^h$, that is nearly the optimal over $\CV^h$ in the graph norm,
$\|\bfe^h\|_{\CG}^2:= \frac{1}{2}\|\bfe^h\|^2 + \|L\bfe^h\|^2$.
The FOSLS and intermediate terms in the Hybrid functional provide a very effective a
posteriori error measure.

We show that the hybrid functional is coercive and continuous in the
graph-like
norm with modest constants, $c_0 = 1/3$ and $c_1=3$; that both $\|\bfe^h \|$ and $\|L
\bfe^h\|$
converge with rates based on standard interpolation bounds; and that if $LL^*$ has full
$H^2$ regularity,
the $\BL^2$ error, $\|\bfe^h\|$, converges with a full power of the discretization
parameter, $h$, faster
than the functional norm. Letting $\tilde{\bu}^h$ denote the optimum over $\CV^h$ in the
graph norm, we also
show that if superposition is used, then $\| \bu^h -\tilde{\bu}^h\|_{\CG}$
converges two powers of $h$ faster than the functional norm. Numerical tests are provided
to confirm the efficiency of the Hybrid method.


\end{document}
